\section*{ТЕРМИНЫ, ОПРЕДЕЛЕНИЯ И СОКРАЩЕНИЯ}
\addcontentsline{toc}{section}{Термины, определения и сокращения}

\subsection*{Термины и определения}

\textbf{Актуарные модели} --- математические и статистические инструменты, используемые для оценки рисков, прогнозирования финансовых последствий неопределённостей, а также для расчёта страховых взносов, управления активами и обязательствами в различных сферах, прежде всего в страховании.

\textbf{Аддитивный принцип} --- метод расчета ожидаемых убытков, при котором итоговая стоимость формируется как сумма произведений вероятностей различных страховых событий на соответствующие выплаты.

\textbf{Бонус-малус система (БМС)} --- механизм индивидуализации страхового тарифа на основе страховой истории клиента, предусматривающий применение скидок за безаварийную езду и надбавок за наличие страховых выплат; в работе рассматривается как традиционный метод, подлежащий оптимизации средствами машинного обучения.

\textbf{Бутстрэппинг} (Bootstrap) --- статистический метод повторной выборки с возвращением, используемый для оценки точности выборочных оценок; в алгоритме Random Forest применяется для декорреляции отдельных деревьев.

\textbf{Бернуллиевское распределение} --- дискретное распределение вероятностей случайной величины, принимающей два возможных значения (0 и 1); используется в работе для моделирования целевой переменной в задаче бинарной классификации.

\textbf{Бинарная кросс-энтропия} --- функция потерь, используемая в задачах бинарной классификации для оценки расхождения между предсказанными вероятностями и истинными метками классов; минимизация данной функции эквивалентна максимизации правдоподобия данных.

\textbf{Валидационная выборка} --- подмножество данных, используемое для настройки гиперпараметров модели и предотвращения переобучения в процессе обучения; не участвует в непосредственном обучении модели.

\textbf{Важность признаков} --- метрика, оценивающая вклад каждого входного параметра в прогноз модели машинного обучения; используется в работе для интерпретации результатов и выявления ключевых факторов риска.

\textbf{Гетерогенные данные} --- данные, содержащие признаки различных типов (числовые, категориальные, текстовые); в работе алгоритм CatBoost применяется именно для работы с такими данными.

\textbf{Гиперпараметры} --- параметры алгоритма машинного обучения, которые задаются до начала обучения и не изменяются в процессе; в работе подбираются методом сеточного поиска и кросс-валидации.

\textbf{Градиентный бустинг} --- ансамблевый алгоритм машинного обучения, основанный на последовательном построении композиции моделей, где каждая следующая модель обучается минимизировать ошибки предыдущих; в работе используется как основной метод прогнозирования вероятности ДТП.

\textbf{Дисбаланс классов} --- ситуация в задачах классификации, при которой количество объектов разных классов существенно различается; в страховых данных характерен значительный перевес объектов без страховых случаев над объектами со страховыми случаями.

\textbf{Динамическое ценообразование} --- подход к формированию стоимости страхового полиса, при котором тариф адаптируется под индивидуальный риск конкретного клиента в режиме, близком к реальному времени.

\textbf{Инференс} --- процесс применения обученной модели машинного обучения к новым данным для получения прогнозов; в работе осуществляется через метод \texttt{predict\_proba()} класса \texttt{InsuranceRiskModel}.

\textbf{Клиент-серверная архитектура} --- архитектурный паттерн, разделяющий систему на клиентскую часть, интерфейс пользователя, и серверную часть; в работе реализована через взаимодействие браузера и Flask-приложения.

\textbf{Коэффициент корреляции Пирсона} --- статистическая мера линейной зависимости между двумя переменными; в работе используется для оценки значимости связи признаков с целевой переменной.

\textbf{Конвейер обработки данных} --- последовательность этапов обработки информации, включающая сбор, очистку, нормализацию данных, обучение модели и получение прогноза; в работе реализован в виде программного класса \texttt{InsuranceRiskModel}.

\textbf{Кросс-валидация} --- метод оценки обобщающей способности модели, при котором данные многократно разделяются на обучающую и проверочную выборки; в работе применяется для подбора гиперпараметров.

\textbf{Кросс-браузерное тестирование} --- проверка корректности работы веб-приложения в различных браузерах для обеспечения универсальной доступности.

\textbf{Логистическая регрессия} --- статистическая модель, используемая для задач бинарной классификации, основанная на применении сигмоидной функции к линейной комбинации признаков; в работе используется как базовая модель для сравнительного анализа.

\textbf{Маргинальные распределения} --- распределения вероятностей отдельных случайных величин, полученные из совместного распределения путем суммирования по остальным переменным; традиционные мультипликативные модели ошибочно предполагают независимость маргинальных распределений.

\textbf{Марковская цепь} --- стохастический процесс с дискретным временем, в котором будущее состояние зависит только от текущего состояния, но не от предыдущих; система бонус-малус описывается как дискретная марковская цепь с 15 классами.

\textbf{Матрица ошибок} --- таблица, отображающая количество истинно положительных, истинно отрицательных, ложноположительных и ложноотрицательных прогнозов модели; используется для детального анализа качества классификации.

\textbf{Мета-модель} --- модель более высокого уровня, обучаемая на предсказаниях базовых моделей в методе стекинга; позволяет компенсировать индивидуальные слабости отдельных алгоритмов.

\textbf{Метод главных компонент} --- метод снижения размерности данных, находящий ортогональные направления максимальной дисперсии; применяется для визуализации многомерных данных и ускорения обучения моделей.

\textbf{Метрики классификации} --- количественные показатели качества работы модели бинарной классификации; в работе используются Accuracy (точность), Precision (точность положительного класса), Recall (полнота), F1-score (гармоническое среднее) и ROC-AUC (площадь под ROC-кривой).

\textbf{Мультиколлинеарность} --- высокая корреляция между двумя или более признаками в модели регрессии, что может приводить к неустойчивости оценок коэффициентов; в работе контролируется через коэффициент вариации инфляции.

\textbf{Мультипликативная модель} --- модель расчета страхового тарифа, в которой итоговая стоимость определяется произведением базовой ставки на набор коэффициентов; традиционная формула ОСАГО является мультипликативной.

\textbf{Обучение с подкреплением} --- раздел машинного обучения, в котором агент обучается выбирать оптимальные действия путем взаимодействия со средой и получения обратной связи; применяется в задачах динамического ценообразования.

\textbf{Ordered Boosting} --- метод обучения в алгоритме CatBoost, использующий перестановки обучающей выборки для расчёта градиентов, что устраняет смещение предсказаний и повышает обобщающую способность модели.

\textbf{Перекрестное субсидирование} --- ситуация в страховании, при которой низкорисковые клиенты фактически финансируют выплаты за высокорисковых из-за усредненных тарифов; разработанная система направлена на снижение данного эффекта.

\textbf{Переобучение} --- ситуация в машинном обучении, при которой модель чрезмерно точно воспроизводит закономерности обучающей выборки, включая шум и выбросы, что приводит к снижению качества предсказаний на новых данных; для предотвращения используется метод Ordered Boosting в алгоритме CatBoost.

\textbf{Персонализированное ценообразование} --- подход к формированию страхового тарифа, учитывающий индивидуальные характеристики и поведение конкретного клиента, в отличие от усредненных категорийных тарифов.

\textbf{Порог классификации} --- граничное значение вероятности, выше которого наблюдение классифицируется как принадлежащее положительному классу; в работе подобран оптимальный порог $\theta^* = 0.30$.

\textbf{Правдоподобие} --- функция, измеряющая соответствие модели наблюдаемым данным; минимизация бинарной кросс-энтропии эквивалентна максимизации правдоподобия при бернуллиевском распределении целевой переменной.

\textbf{Производные признаки} --- новые признаки, создаваемые на основе исходных переменных для выявления нелинейных взаимодействий и повышения предсказательной способности модели; в работе сгенерировано 10 производных признаков.

\textbf{Ранняя остановка} --- метод регуляризации, при котором обучение модели прекращается, если качество на валидационной выборке не улучшается в течение заданного числа итераций; в работе используется \texttt{early\_stopping\_rounds = 100}.

\textbf{Регуляризация} --- набор методов, направленных на предотвращение переобучения модели путем введения штрафов за сложность; в XGBoost реализуется через штраф $\Omega(f_k)$ за сложность дерева.

\textbf{Сериализация} --- процесс преобразования объекта, в данном случае обученной модели, в формат, пригодный для хранения и последующего восстановления; в работе модель сохраняется в бинарном формате \texttt{.cbm}.

\textbf{Совместное распределение} --- распределение вероятностей системы случайных величин, описывающее их одновременное поведение; традиционные модели ошибочно факторизуют совместное распределение как произведение маргинальных.

\textbf{Стекинг} --- ансамблевый метод, объединяющий предсказания нескольких базовых моделей через мета-обучение для повышения точности прогнозов.

\textbf{Стратификация} --- метод разделения выборки на подмножества с сохранением пропорций классов в каждом подмножестве; в работе применяется для обеспечения баланса классов в обучающей, валидационной и тестовой выборках.

\textbf{Страховой риск} --- предполагаемое событие, обладающее признаками вероятности и случайности, при наступлении которого возникает обязательство страховщика произвести страховую выплату; в контексте работы: вероятность наступления дорожно-транспортного происшествия.

\textbf{Страховой тариф} --- цена страхового риска, выраженная в денежной сумме, относящаяся к единице страховой суммы или процентной доле от нее; в работе рассчитывается на основе прогнозной вероятности ДТП, полученной моделью машинного обучения.

\textbf{Target-based Encoding} --- метод кодирования категориальных признаков, при котором категория заменяется на статистически сглаженную оценку условной вероятности целевой переменной; реализован в алгоритме CatBoost.

\textbf{Телематические данные} --- информация о поведении водителя и состоянии транспортного средства, собираемая специальными устройствами, такими как скорость, ускорение, резкость торможения, геолокация; в работе используется датасет \textit{Levin Vehicle Telematics}.

\textbf{Телематическое страхование} --- вид страхования, при котором страховой тариф формируется на основе данных о реальном поведении водителя, такими как скорость, пробег, стиль вождения, фиксируемых специальными устройствами; рассматривается в работе как перспективное направление расширения признаков модели.

\textbf{Тестовая выборка} --- подмножество данных, не участвующее в обучении и настройке модели, используемое для финальной оценки качества модели на новых данных.

\textbf{Функция активации} --- нелинейное преобразование, применяемое к взвешенной сумме входов нейрона в искусственной нейронной сети; обеспечивает способность сети аппроксимировать сложные нелинейные зависимости.

\textbf{Функция потерь} --- математическая функция, измеряющая расхождение между предсказаниями модели и истинными значениями; в работе используется функция бинарной кросс-энтропии.

\textbf{Целевая переменная} --- бинарный индикатор факта дорожно-транспортного происшествия, соответствующий значению 1 при наличии страхового случая и 0 при его отсутствии; является объектом прогнозирования в задаче классификации.

\textbf{Категориальный признак} --- переменная, которая принимает дискретный набор значений; в работе обрабатывается алгоритмом CatBoost без предварительного кодирования.

\textbf{Базовая ставка тарифа} --- минимальная стоимость страхового полиса, устанавливаемая страховой компанией для низкорисковых клиентов; служит основой для динамического расчета итогового тарифа.

\textbf{Коэффициент чувствительности к риску} --- параметр, отвечающий за степень влияния прогнозируемой вероятности ДТП на итоговую стоимость страхового полиса; регулируется экспертно в зависимости от бизнес-политики страховщика.

\textbf{Адаптивная верстка} --- подход к проектированию веб-интерфейса, при котором элементы автоматически подстраиваются под размеры и ориентацию экрана устройства пользователя.

\textbf{Маршрутизация} --- механизм сопоставления URL-адресов с функциями-обработчиками в веб-приложении; в работе реализована через систему декораторов Flask.

\textbf{Бортовая диагностика} --- система самодиагностики транспортного средства, фиксирующая коды неисправностей и технические параметры; данные OBD используются в работе как телематические признаки.

\textbf{Коды диагностических неисправностей} (DTC, Diagnostic Trouble Codes) --- стандартизированные коды ошибок, генерируемые бортовой системой автомобиля при обнаружении технических неисправностей; загрузка DTC-файлов реализована в веб-приложении.

\textbf{Коэффициент бонус-малус} (КБМ) --- показатель, применяемый в системе ОСАГО для индивидуализации тарифа на основе истории страховых случаев водителя; в работе рассчитывается по формуле динамического ценообразования.

\textbf{Случайный лес} (Random Forest) --- ансамблевый алгоритм машинного обучения, состоящий из множества решающих деревьев, обученных на бутстрэп-выборках со случайным подмножеством признаков; итоговый прогноз формируется усреднением.

\textbf{Обучающая выборка} --- подмножество данных, используемое непосредственно для обучения модели машинного обучения путем оптимизации функции потерь.

\textbf{Диаграмма прецедентов} --- диаграмма, отображающая основные сценарии взаимодействия пользователя с системой; используется в работе для описания функциональных требований.

\textbf{Диаграмма пользовательского сценария} --- диаграмма, визуализирующая последовательность действий пользователя при взаимодействии с системой от начала до получения результата.

\textbf{Граничные значения} --- экстремальные допустимые значения параметров, используемые при тестировании для проверки корректности обработки системой крайних случаев.

\textbf{Верификация} --- процесс проверки соответствия разработанной системы заданным требованиям и спецификациям; в работе осуществляется на всех этапах разработки модели.

\textbf{Клиппинг выбросов} --- метод обработки статистических аномалий, при котором значения признаков ограничиваются физиологически и практически обоснованными диапазонами.

\textbf{Сеточный поиск} --- метод подбора гиперпараметров, при котором производится перебор всех возможных комбинаций значений из заданного множества; в работе применяется совместно с кросс-валидацией.

\textbf{Совместное распределение} --- распределение вероятностей системы случайных величин, описывающее их одновременное поведение; традиционные модели ошибочно факторизуют совместное распределение как произведение маргинальных.

\subsection*{Перечень сокращений}

{\footnotesize
\begin{tabular}{@{}lp{10cm}@{}}
ОСАГО & Обязательное страхование гражданской ответственности владельцев транспортных средств \\[4pt]
КАСКО & Комплексное автострахование, кроме ответственности (добровольное страхование транспортного средства от ущерба и угона) \\[4pt]
ДТП & Дорожно-транспортное происшествие \\[4pt]
ПДД & Правила дорожного движения \\[4pt]
ЦБ РФ & Центральный банк Российской Федерации (регулятор страхового рынка) \\[4pt]
КБМ & Коэффициент бонус-малус \\[4pt]
МЛ / ML & Machine Learning --- машинное обучение \\[4pt]
ИИ / AI & Artificial Intelligence --- искусственный интеллект \\[4pt]
EDA & Exploratory Data Analysis --- разведочный анализ данных \\[4pt]
CatBoost & Categorical Boosting --- алгоритм градиентного бустинга \\[4pt]
XGBoost & eXtreme Gradient Boosting --- алгоритм градиентного бустинга \\[4pt]
LightGBM & Light Gradient Boosting Machine --- алгоритм градиентного бустинга \\[4pt]
ROC-AUC & Receiver Operating Characteristic -- Area Under Curve \\[4pt]
F1-score & Гармоническое среднее между Precision и Recall \\[4pt]
Precision & Доля правильно предсказанных положительных случаев \\[4pt]
Recall & Доля правильно предсказанных случаев среди всех истинных \\[4pt]
Accuracy & Доля правильно классифицированных объектов \\[4pt]
UBI & Usage-Based Insurance --- телематическое страхование \\[4pt]
OBD & On-Board Diagnostics --- бортовая диагностика \\[4pt]
DTC & Diagnostic Trouble Codes --- коды диагностических неисправностей \\[4pt]
VIF & Variance Inflation Factor --- коэффициент вариации инфляции \\[4pt]
PCA & Principal Component Analysis --- метод главных компонент \\[4pt]
MDP & Markov Decision Process --- марковский процесс принятия решений \\[4pt]
MVC & Model-View-Controller --- архитектурный паттерн «Модель-Представление-Контроллер» \\[4pt]
\end{tabular}

\newpage

\begin{tabular}{@{}lp{10cm}@{}}
REST & Representational State Transfer --- архитектурный стиль взаимодействия компонентов \\[4pt]
API & Application Programming Interface --- программный интерфейс приложения \\[4pt]
HTTP & HyperText Transfer Protocol --- протокол передачи гипертекста \\[4pt]
JSON & JavaScript Object Notation --- формат обмена данными \\[4pt]
CSV & Comma-Separated Values --- формат представления табличных данных \\[4pt]
PDF & Portable Document Format --- формат электронных документов \\[4pt]
PNG & Portable Network Graphics --- формат растровых изображений \\[4pt]
HTML & HyperText Markup Language --- язык разметки гипертекста \\[4pt]
CSS & Cascading Style Sheets --- каскадные таблицы стилей \\[4pt]
UX & User Experience --- пользовательский опыт \\[4pt]
UI & User Interface --- пользовательский интерфейс \\[4pt]
\end{tabular}
}